\section{Prinsip Kerja Setiap Komponen}
\label{sec:prinsip-komponen}

Bagian ini merupakan fokus utama FGD. Setiap komponen dibahas secara mandiri
agar pembaca memahami cara kerja, karakteristik sinyal, kebutuhan antarmuka,
dan batasan komponen sebelum melihat penerapannya dalam proyek.

\subsection{Cara Menggunakan Template Komponen}
Setiap komponen ditempatkan pada file terpisah di direktori
\texttt{sections/components/}. Gunakan struktur yang sama untuk setiap
komponen agar hasil eksplorasi dapat dibandingkan secara sistematis. Isi
bagian yang tidak relevan dengan \textit{Tidak berlaku} dan tambahkan
subsubsection khusus hanya jika memang diperlukan oleh karakteristik komponen.

% =====================================================================
% TEMPLATE ELABORASI SATU KOMPONEN
% Salin file ini menjadi komponen_NAMA.tex, lalu ganti placeholder.
% Jangan meng-input file template ini secara langsung ke bab.
% =====================================================================

\subsection{Komponen 1}
\label{subsec:komponen-01}

% ---------------------------------------------------------------------
\subsubsection{Identitas dan Fungsi}
Jelaskan nama komponen, produsen/model, kategori komponen (sensor, aktuator,
modul, atau komponen pendukung), fungsi utama, serta alasan komponen dipilih
untuk dieksplorasi.

\begin{table}[H]
\centering
\caption{Identitas komponen}
\begin{tabularx}{\textwidth}{@{}lX@{}}
\toprule
Parameter & Keterangan \\
\midrule
Nama komponen & \dots \\
Model/part number & \dots \\
Jenis & Sensor/Aktuator/Modul/Komponen pendukung \\
Fungsi utama & \dots \\
Sumber dokumentasi & \dots \\
\bottomrule
\end{tabularx}
\end{table}

% ---------------------------------------------------------------------
\subsubsection{Prinsip Kerja}
Jelaskan fenomena fisika atau elektronika yang membuat komponen memberikan
respons tertentu. Susun penjelasan dari kondisi masukan, proses internal,
hingga keluaran. Jika relevan, sertakan diagram, kurva karakteristik, atau
persamaan matematis.

% Contoh format persamaan:
% \begin{equation}
%   y = f(x)
%   \label{eq:komponen-nama}
% \end{equation}

\subsubsection{Alur Kerja Internal}
Uraikan alur \textbf{input \textrightarrow proses \textrightarrow output}.
Untuk sensor, jelaskan perubahan besaran fisik menjadi sinyal listrik.
Untuk aktuator, jelaskan perubahan sinyal listrik menjadi aksi fisik.
Jika komponen berupa modul digital, jelaskan proses pengolahan dan komunikasi
data secara konseptual.

% ---------------------------------------------------------------------
\subsubsection{Besaran yang Diukur atau Dikendalikan}
Jelaskan besaran utama, satuan, rentang pengukuran atau kendali, resolusi,
sensitivitas, akurasi, serta parameter lain yang relevan.

\begin{table}[H]
\centering
\caption{Karakteristik besaran komponen}
\begin{tabularx}{\textwidth}{@{}lXX@{}}
\toprule
Parameter & Nilai & Keterangan \\
\midrule
Besaran utama & \dots & \dots \\
Rentang & \dots & \dots \\
Resolusi & \dots & \dots \\
Akurasi/sensitivitas & \dots & \dots \\
\bottomrule
\end{tabularx}
\end{table}

% ---------------------------------------------------------------------
\subsubsection{Karakteristik Sinyal}
Identifikasi jenis sinyal dan jelaskan bentuk serta karakteristiknya. Bahas
apakah sinyal bersifat analog, digital, PWM, UART, I\textsuperscript{2}C,
SPI, atau protokol lain. Jelaskan level logika, rentang tegangan, frekuensi,
kecepatan data, resolusi ADC/DAC, atau parameter lain yang diperlukan.

\begin{table}[H]
\centering
\caption{Karakteristik sinyal}
\begin{tabularx}{\textwidth}{@{}lXX@{}}
\toprule
Aspek & Nilai & Keterangan \\
\midrule
Jenis sinyal & \dots & \dots \\
Level logika & \dots & \dots \\
Rentang tegangan & \dots & \dots \\
Frekuensi/baud rate & \dots & \dots \\
\bottomrule
\end{tabularx}
\end{table}

% ---------------------------------------------------------------------
\subsubsection{Pin, Catu Daya, dan Antarmuka}
Jelaskan fungsi setiap pin dan kebutuhan catu daya. Jika terdapat lebih dari
satu mode antarmuka, jelaskan perbedaannya dan alasan pemilihannya.

\begin{table}[H]
\centering
\caption{Pin dan antarmuka komponen}
\begin{tabularx}{\textwidth}{@{}lXX@{}}
\toprule
Pin/Antarmuka & Fungsi & Keterangan \\
\midrule
VCC/VIN & Catu daya & \dots \\
GND & Referensi & \dots \\
OUT/SDA/SCL/dll. & Sinyal & \dots \\
\bottomrule
\end{tabularx}
\end{table}

\subsubsection{Kebutuhan Catu Daya}
Jelaskan tegangan kerja, konsumsi arus, kondisi minimum/maksimum, kebutuhan
regulator, decoupling, serta kompatibilitas level tegangan dengan
mikrokontroler.

% ---------------------------------------------------------------------
\subsubsection{Spesifikasi Penting dari Datasheet}
Ambil hanya parameter yang relevan terhadap pemakaian komponen. Jangan sekadar
menyalin seluruh tabel datasheet; jelaskan arti setiap parameter dan mengapa
parameter tersebut penting untuk eksperimen.

\begin{table}[H]
\centering
\caption{Spesifikasi penting komponen}
\begin{tabular}{@{}lcc@{}}
\toprule
Parameter & Nilai & Satuan \\
\midrule
Tegangan kerja & \dots & \si{\volt} \\
Rentang ukur/keluaran & \dots & \dots \\
Resolusi/akurasi & \dots & \dots \\
Waktu respons & \dots & \si{\milli\second} \\
Arus konsumsi & \dots & \si{\milli\ampere} \\
\bottomrule
\end{tabular}
\end{table}

Sumber utama parameter: \cite{datasheet}.

% ---------------------------------------------------------------------
\subsubsection{Hubungan dengan Mikrokontroler}
Jelaskan fitur mikrokontroler yang digunakan, misalnya GPIO, ADC, timer,
interrupt, PWM, UART, I\textsuperscript{2}C, atau SPI. Kaitkan kebutuhan
tersebut dengan karakteristik sinyal komponen dan konsep pada Bagian B.
Referensi mikrokontroler dapat mengacu pada~\cite{mcu}.

\begin{table}[H]
\centering
\caption{Pemetaan komponen ke mikrokontroler}
\begin{tabularx}{\textwidth}{@{}lXX@{}}
\toprule
Kebutuhan komponen & Fitur MCU & Alasan pemilihan \\
\midrule
Input/output & \dots & \dots \\
Antarmuka komunikasi & \dots & \dots \\
Pemrosesan sinyal & \dots & \dots \\
\bottomrule
\end{tabularx}
\end{table}

% ---------------------------------------------------------------------
\subsubsection{Rangkaian Antarmuka}
Tampilkan skema koneksi komponen dengan mikrokontroler. Jelaskan aliran
sinyal dan fungsi setiap komponen eksternal. Jika ada resistor, kapasitor,
transistor, regulator, atau level shifter, sertakan dasar pemilihannya.

\begin{figure}[H]
\centering
% Ganti dengan gambar rangkaian aktual, misalnya:
% \includegraphics[width=0.85\textwidth]{images/rangkaian-nama-komponen.png}
\fbox{\parbox[c][4cm][c]{0.85\textwidth}{\centering Tempatkan skema rangkaian di sini}}
\caption{Rangkaian antarmuka Komponen 1 dengan mikrokontroler}
\label{fig:rangkaian-komponen-01}
\end{figure}

\subsubsection{Perhitungan dan Justifikasi Rangkaian}
Tunjukkan perhitungan yang diperlukan untuk menentukan nilai komponen
pendukung atau mengubah besaran sinyal. Jelaskan asumsi yang digunakan dan
bandingkan hasil perhitungan dengan batas datasheet.

% ---------------------------------------------------------------------
\subsubsection{Contoh Penggunaan Sederhana}
Buat eksperimen kecil yang hanya bertujuan memahami komponen. Jelaskan
input, prosedur, output yang diharapkan, dan indikator bahwa komponen bekerja
sesuai prinsipnya. Eksperimen ini harus dapat berdiri sendiri sebelum masuk
ke proyek utama.

\subsubsection{Program atau Algoritma Dasar}
Tampilkan bagian kode minimum yang diperlukan untuk membaca atau mengendalikan
komponen. Jelaskan fungsi bagian penting kode, bukan hanya menampilkan kode.

\begin{lstlisting}[caption={Contoh program dasar komponen}]
// Tulis kode minimum untuk menguji komponen di sini.
\end{lstlisting}

% ---------------------------------------------------------------------
\subsubsection{Hasil Eksplorasi}
Dokumentasikan hasil pengamatan, nilai pembacaan, respons komponen, foto,
grafik, atau bukti lain. Bedakan antara hasil yang diharapkan berdasarkan
teori dan hasil aktual.

\begin{table}[H]
\centering
\caption{Contoh format hasil eksplorasi}
\begin{tabularx}{\textwidth}{@{}lXX@{}}
\toprule
Kondisi/Input & Hasil Teoretis & Hasil Pengamatan \\
\midrule
Kondisi 1 & \dots & \dots \\
Kondisi 2 & \dots & \dots \\
Kondisi 3 & \dots & \dots \\
\bottomrule
\end{tabularx}
\end{table}

% ---------------------------------------------------------------------
\subsubsection{Analisis dan Interpretasi}
Hubungkan hasil eksplorasi dengan prinsip kerja dan spesifikasi datasheet.
Jelaskan apakah hasil sesuai teori, penyebab perbedaan, serta implikasinya
terhadap penggunaan komponen pada proyek.

\subsubsection{Keterbatasan dan Sumber Kesalahan}
Bahas noise, toleransi komponen, kesalahan pembacaan, kalibrasi, temperatur,
kondisi lingkungan, keterbatasan library, keterbatasan catu daya, batas
komunikasi, atau faktor lain yang dapat memengaruhi hasil.

\subsubsection{Kesimpulan Pemahaman Komponen}
Tuliskan secara ringkas apa yang dapat disimpulkan mengenai prinsip kerja,
karakteristik sinyal, kebutuhan mikrokontroler, dan kondisi penggunaan
komponen. Tekankan bukti yang menunjukkan bahwa prinsip kerja telah dipahami.

% =====================================================================
% TEMPLATE ELABORASI SATU KOMPONEN
% Salin file ini menjadi komponen_NAMA.tex, lalu ganti placeholder.
% Jangan meng-input file template ini secara langsung ke bab.
% =====================================================================

\subsection{Komponen 2}
\label{subsec:komponen-02}

% ---------------------------------------------------------------------
\subsubsection{Identitas dan Fungsi}
Jelaskan nama komponen, produsen/model, kategori komponen (sensor, aktuator,
modul, atau komponen pendukung), fungsi utama, serta alasan komponen dipilih
untuk dieksplorasi.

\begin{table}[H]
\centering
\caption{Identitas komponen}
\begin{tabularx}{\textwidth}{@{}lX@{}}
\toprule
Parameter & Keterangan \\
\midrule
Nama komponen & \dots \\
Model/part number & \dots \\
Jenis & Sensor/Aktuator/Modul/Komponen pendukung \\
Fungsi utama & \dots \\
Sumber dokumentasi & \dots \\
\bottomrule
\end{tabularx}
\end{table}

% ---------------------------------------------------------------------
\subsubsection{Prinsip Kerja}
Jelaskan fenomena fisika atau elektronika yang membuat komponen memberikan
respons tertentu. Susun penjelasan dari kondisi masukan, proses internal,
hingga keluaran. Jika relevan, sertakan diagram, kurva karakteristik, atau
persamaan matematis.

% Contoh format persamaan:
% \begin{equation}
%   y = f(x)
%   \label{eq:komponen-nama}
% \end{equation}

\subsubsection{Alur Kerja Internal}
Uraikan alur \textbf{input \textrightarrow proses \textrightarrow output}.
Untuk sensor, jelaskan perubahan besaran fisik menjadi sinyal listrik.
Untuk aktuator, jelaskan perubahan sinyal listrik menjadi aksi fisik.
Jika komponen berupa modul digital, jelaskan proses pengolahan dan komunikasi
data secara konseptual.

% ---------------------------------------------------------------------
\subsubsection{Besaran yang Diukur atau Dikendalikan}
Jelaskan besaran utama, satuan, rentang pengukuran atau kendali, resolusi,
sensitivitas, akurasi, serta parameter lain yang relevan.

\begin{table}[H]
\centering
\caption{Karakteristik besaran komponen}
\begin{tabularx}{\textwidth}{@{}lXX@{}}
\toprule
Parameter & Nilai & Keterangan \\
\midrule
Besaran utama & \dots & \dots \\
Rentang & \dots & \dots \\
Resolusi & \dots & \dots \\
Akurasi/sensitivitas & \dots & \dots \\
\bottomrule
\end{tabularx}
\end{table}

% ---------------------------------------------------------------------
\subsubsection{Karakteristik Sinyal}
Identifikasi jenis sinyal dan jelaskan bentuk serta karakteristiknya. Bahas
apakah sinyal bersifat analog, digital, PWM, UART, I\textsuperscript{2}C,
SPI, atau protokol lain. Jelaskan level logika, rentang tegangan, frekuensi,
kecepatan data, resolusi ADC/DAC, atau parameter lain yang diperlukan.

\begin{table}[H]
\centering
\caption{Karakteristik sinyal}
\begin{tabularx}{\textwidth}{@{}lXX@{}}
\toprule
Aspek & Nilai & Keterangan \\
\midrule
Jenis sinyal & \dots & \dots \\
Level logika & \dots & \dots \\
Rentang tegangan & \dots & \dots \\
Frekuensi/baud rate & \dots & \dots \\
\bottomrule
\end{tabularx}
\end{table}

% ---------------------------------------------------------------------
\subsubsection{Pin, Catu Daya, dan Antarmuka}
Jelaskan fungsi setiap pin dan kebutuhan catu daya. Jika terdapat lebih dari
satu mode antarmuka, jelaskan perbedaannya dan alasan pemilihannya.

\begin{table}[H]
\centering
\caption{Pin dan antarmuka komponen}
\begin{tabularx}{\textwidth}{@{}lXX@{}}
\toprule
Pin/Antarmuka & Fungsi & Keterangan \\
\midrule
VCC/VIN & Catu daya & \dots \\
GND & Referensi & \dots \\
OUT/SDA/SCL/dll. & Sinyal & \dots \\
\bottomrule
\end{tabularx}
\end{table}

\subsubsection{Kebutuhan Catu Daya}
Jelaskan tegangan kerja, konsumsi arus, kondisi minimum/maksimum, kebutuhan
regulator, decoupling, serta kompatibilitas level tegangan dengan
mikrokontroler.

% ---------------------------------------------------------------------
\subsubsection{Spesifikasi Penting dari Datasheet}
Ambil hanya parameter yang relevan terhadap pemakaian komponen. Jangan sekadar
menyalin seluruh tabel datasheet; jelaskan arti setiap parameter dan mengapa
parameter tersebut penting untuk eksperimen.

\begin{table}[H]
\centering
\caption{Spesifikasi penting komponen}
\begin{tabular}{@{}lcc@{}}
\toprule
Parameter & Nilai & Satuan \\
\midrule
Tegangan kerja & \dots & \si{\volt} \\
Rentang ukur/keluaran & \dots & \dots \\
Resolusi/akurasi & \dots & \dots \\
Waktu respons & \dots & \si{\milli\second} \\
Arus konsumsi & \dots & \si{\milli\ampere} \\
\bottomrule
\end{tabular}
\end{table}

Sumber utama parameter: \cite{datasheet}.

% ---------------------------------------------------------------------
\subsubsection{Hubungan dengan Mikrokontroler}
Jelaskan fitur mikrokontroler yang digunakan, misalnya GPIO, ADC, timer,
interrupt, PWM, UART, I\textsuperscript{2}C, atau SPI. Kaitkan kebutuhan
tersebut dengan karakteristik sinyal komponen dan konsep pada Bagian B.
Referensi mikrokontroler dapat mengacu pada~\cite{mcu}.

\begin{table}[H]
\centering
\caption{Pemetaan komponen ke mikrokontroler}
\begin{tabularx}{\textwidth}{@{}lXX@{}}
\toprule
Kebutuhan komponen & Fitur MCU & Alasan pemilihan \\
\midrule
Input/output & \dots & \dots \\
Antarmuka komunikasi & \dots & \dots \\
Pemrosesan sinyal & \dots & \dots \\
\bottomrule
\end{tabularx}
\end{table}

% ---------------------------------------------------------------------
\subsubsection{Rangkaian Antarmuka}
Tampilkan skema koneksi komponen dengan mikrokontroler. Jelaskan aliran
sinyal dan fungsi setiap komponen eksternal. Jika ada resistor, kapasitor,
transistor, regulator, atau level shifter, sertakan dasar pemilihannya.

\begin{figure}[H]
\centering
% Ganti dengan gambar rangkaian aktual, misalnya:
% \includegraphics[width=0.85\textwidth]{images/rangkaian-nama-komponen.png}
\fbox{\parbox[c][4cm][c]{0.85\textwidth}{\centering Tempatkan skema rangkaian di sini}}
\caption{Rangkaian antarmuka Komponen 2 dengan mikrokontroler}
\label{fig:rangkaian-komponen-02}
\end{figure}

\subsubsection{Perhitungan dan Justifikasi Rangkaian}
Tunjukkan perhitungan yang diperlukan untuk menentukan nilai komponen
pendukung atau mengubah besaran sinyal. Jelaskan asumsi yang digunakan dan
bandingkan hasil perhitungan dengan batas datasheet.

% ---------------------------------------------------------------------
\subsubsection{Contoh Penggunaan Sederhana}
Buat eksperimen kecil yang hanya bertujuan memahami komponen. Jelaskan
input, prosedur, output yang diharapkan, dan indikator bahwa komponen bekerja
sesuai prinsipnya. Eksperimen ini harus dapat berdiri sendiri sebelum masuk
ke proyek utama.

\subsubsection{Program atau Algoritma Dasar}
Tampilkan bagian kode minimum yang diperlukan untuk membaca atau mengendalikan
komponen. Jelaskan fungsi bagian penting kode, bukan hanya menampilkan kode.

\begin{lstlisting}[caption={Contoh program dasar komponen}]
// Tulis kode minimum untuk menguji komponen di sini.
\end{lstlisting}

% ---------------------------------------------------------------------
\subsubsection{Hasil Eksplorasi}
Dokumentasikan hasil pengamatan, nilai pembacaan, respons komponen, foto,
grafik, atau bukti lain. Bedakan antara hasil yang diharapkan berdasarkan
teori dan hasil aktual.

\begin{table}[H]
\centering
\caption{Contoh format hasil eksplorasi}
\begin{tabularx}{\textwidth}{@{}lXX@{}}
\toprule
Kondisi/Input & Hasil Teoretis & Hasil Pengamatan \\
\midrule
Kondisi 1 & \dots & \dots \\
Kondisi 2 & \dots & \dots \\
Kondisi 3 & \dots & \dots \\
\bottomrule
\end{tabularx}
\end{table}

% ---------------------------------------------------------------------
\subsubsection{Analisis dan Interpretasi}
Hubungkan hasil eksplorasi dengan prinsip kerja dan spesifikasi datasheet.
Jelaskan apakah hasil sesuai teori, penyebab perbedaan, serta implikasinya
terhadap penggunaan komponen pada proyek.

\subsubsection{Keterbatasan dan Sumber Kesalahan}
Bahas noise, toleransi komponen, kesalahan pembacaan, kalibrasi, temperatur,
kondisi lingkungan, keterbatasan library, keterbatasan catu daya, batas
komunikasi, atau faktor lain yang dapat memengaruhi hasil.

\subsubsection{Kesimpulan Pemahaman Komponen}
Tuliskan secara ringkas apa yang dapat disimpulkan mengenai prinsip kerja,
karakteristik sinyal, kebutuhan mikrokontroler, dan kondisi penggunaan
komponen. Tekankan bukti yang menunjukkan bahwa prinsip kerja telah dipahami.

% =====================================================================
% TEMPLATE ELABORASI SATU KOMPONEN
% Salin file ini menjadi komponen_NAMA.tex, lalu ganti placeholder.
% Jangan meng-input file template ini secara langsung ke bab.
% =====================================================================

\subsection{Komponen 3}
\label{subsec:komponen-03}

% ---------------------------------------------------------------------
\subsubsection{Identitas dan Fungsi}
Jelaskan nama komponen, produsen/model, kategori komponen (sensor, aktuator,
modul, atau komponen pendukung), fungsi utama, serta alasan komponen dipilih
untuk dieksplorasi.

\begin{table}[H]
\centering
\caption{Identitas komponen}
\begin{tabularx}{\textwidth}{@{}lX@{}}
\toprule
Parameter & Keterangan \\
\midrule
Nama komponen & \dots \\
Model/part number & \dots \\
Jenis & Sensor/Aktuator/Modul/Komponen pendukung \\
Fungsi utama & \dots \\
Sumber dokumentasi & \dots \\
\bottomrule
\end{tabularx}
\end{table}

% ---------------------------------------------------------------------
\subsubsection{Prinsip Kerja}
Jelaskan fenomena fisika atau elektronika yang membuat komponen memberikan
respons tertentu. Susun penjelasan dari kondisi masukan, proses internal,
hingga keluaran. Jika relevan, sertakan diagram, kurva karakteristik, atau
persamaan matematis.

% Contoh format persamaan:
% \begin{equation}
%   y = f(x)
%   \label{eq:komponen-nama}
% \end{equation}

\subsubsection{Alur Kerja Internal}
Uraikan alur \textbf{input \textrightarrow proses \textrightarrow output}.
Untuk sensor, jelaskan perubahan besaran fisik menjadi sinyal listrik.
Untuk aktuator, jelaskan perubahan sinyal listrik menjadi aksi fisik.
Jika komponen berupa modul digital, jelaskan proses pengolahan dan komunikasi
data secara konseptual.

% ---------------------------------------------------------------------
\subsubsection{Besaran yang Diukur atau Dikendalikan}
Jelaskan besaran utama, satuan, rentang pengukuran atau kendali, resolusi,
sensitivitas, akurasi, serta parameter lain yang relevan.

\begin{table}[H]
\centering
\caption{Karakteristik besaran komponen}
\begin{tabularx}{\textwidth}{@{}lXX@{}}
\toprule
Parameter & Nilai & Keterangan \\
\midrule
Besaran utama & \dots & \dots \\
Rentang & \dots & \dots \\
Resolusi & \dots & \dots \\
Akurasi/sensitivitas & \dots & \dots \\
\bottomrule
\end{tabularx}
\end{table}

% ---------------------------------------------------------------------
\subsubsection{Karakteristik Sinyal}
Identifikasi jenis sinyal dan jelaskan bentuk serta karakteristiknya. Bahas
apakah sinyal bersifat analog, digital, PWM, UART, I\textsuperscript{2}C,
SPI, atau protokol lain. Jelaskan level logika, rentang tegangan, frekuensi,
kecepatan data, resolusi ADC/DAC, atau parameter lain yang diperlukan.

\begin{table}[H]
\centering
\caption{Karakteristik sinyal}
\begin{tabularx}{\textwidth}{@{}lXX@{}}
\toprule
Aspek & Nilai & Keterangan \\
\midrule
Jenis sinyal & \dots & \dots \\
Level logika & \dots & \dots \\
Rentang tegangan & \dots & \dots \\
Frekuensi/baud rate & \dots & \dots \\
\bottomrule
\end{tabularx}
\end{table}

% ---------------------------------------------------------------------
\subsubsection{Pin, Catu Daya, dan Antarmuka}
Jelaskan fungsi setiap pin dan kebutuhan catu daya. Jika terdapat lebih dari
satu mode antarmuka, jelaskan perbedaannya dan alasan pemilihannya.

\begin{table}[H]
\centering
\caption{Pin dan antarmuka komponen}
\begin{tabularx}{\textwidth}{@{}lXX@{}}
\toprule
Pin/Antarmuka & Fungsi & Keterangan \\
\midrule
VCC/VIN & Catu daya & \dots \\
GND & Referensi & \dots \\
OUT/SDA/SCL/dll. & Sinyal & \dots \\
\bottomrule
\end{tabularx}
\end{table}

\subsubsection{Kebutuhan Catu Daya}
Jelaskan tegangan kerja, konsumsi arus, kondisi minimum/maksimum, kebutuhan
regulator, decoupling, serta kompatibilitas level tegangan dengan
mikrokontroler.

% ---------------------------------------------------------------------
\subsubsection{Spesifikasi Penting dari Datasheet}
Ambil hanya parameter yang relevan terhadap pemakaian komponen. Jangan sekadar
menyalin seluruh tabel datasheet; jelaskan arti setiap parameter dan mengapa
parameter tersebut penting untuk eksperimen.

\begin{table}[H]
\centering
\caption{Spesifikasi penting komponen}
\begin{tabular}{@{}lcc@{}}
\toprule
Parameter & Nilai & Satuan \\
\midrule
Tegangan kerja & \dots & \si{\volt} \\
Rentang ukur/keluaran & \dots & \dots \\
Resolusi/akurasi & \dots & \dots \\
Waktu respons & \dots & \si{\milli\second} \\
Arus konsumsi & \dots & \si{\milli\ampere} \\
\bottomrule
\end{tabular}
\end{table}

Sumber utama parameter: \cite{datasheet}.

% ---------------------------------------------------------------------
\subsubsection{Hubungan dengan Mikrokontroler}
Jelaskan fitur mikrokontroler yang digunakan, misalnya GPIO, ADC, timer,
interrupt, PWM, UART, I\textsuperscript{2}C, atau SPI. Kaitkan kebutuhan
tersebut dengan karakteristik sinyal komponen dan konsep pada Bagian B.
Referensi mikrokontroler dapat mengacu pada~\cite{mcu}.

\begin{table}[H]
\centering
\caption{Pemetaan komponen ke mikrokontroler}
\begin{tabularx}{\textwidth}{@{}lXX@{}}
\toprule
Kebutuhan komponen & Fitur MCU & Alasan pemilihan \\
\midrule
Input/output & \dots & \dots \\
Antarmuka komunikasi & \dots & \dots \\
Pemrosesan sinyal & \dots & \dots \\
\bottomrule
\end{tabularx}
\end{table}

% ---------------------------------------------------------------------
\subsubsection{Rangkaian Antarmuka}
Tampilkan skema koneksi komponen dengan mikrokontroler. Jelaskan aliran
sinyal dan fungsi setiap komponen eksternal. Jika ada resistor, kapasitor,
transistor, regulator, atau level shifter, sertakan dasar pemilihannya.

\begin{figure}[H]
\centering
% Ganti dengan gambar rangkaian aktual, misalnya:
% \includegraphics[width=0.85\textwidth]{images/rangkaian-nama-komponen.png}
\fbox{\parbox[c][4cm][c]{0.85\textwidth}{\centering Tempatkan skema rangkaian di sini}}
\caption{Rangkaian antarmuka Komponen 3 dengan mikrokontroler}
\label{fig:rangkaian-komponen-03}
\end{figure}

\subsubsection{Perhitungan dan Justifikasi Rangkaian}
Tunjukkan perhitungan yang diperlukan untuk menentukan nilai komponen
pendukung atau mengubah besaran sinyal. Jelaskan asumsi yang digunakan dan
bandingkan hasil perhitungan dengan batas datasheet.

% ---------------------------------------------------------------------
\subsubsection{Contoh Penggunaan Sederhana}
Buat eksperimen kecil yang hanya bertujuan memahami komponen. Jelaskan
input, prosedur, output yang diharapkan, dan indikator bahwa komponen bekerja
sesuai prinsipnya. Eksperimen ini harus dapat berdiri sendiri sebelum masuk
ke proyek utama.

\subsubsection{Program atau Algoritma Dasar}
Tampilkan bagian kode minimum yang diperlukan untuk membaca atau mengendalikan
komponen. Jelaskan fungsi bagian penting kode, bukan hanya menampilkan kode.

\begin{lstlisting}[caption={Contoh program dasar komponen}]
// Tulis kode minimum untuk menguji komponen di sini.
\end{lstlisting}

% ---------------------------------------------------------------------
\subsubsection{Hasil Eksplorasi}
Dokumentasikan hasil pengamatan, nilai pembacaan, respons komponen, foto,
grafik, atau bukti lain. Bedakan antara hasil yang diharapkan berdasarkan
teori dan hasil aktual.

\begin{table}[H]
\centering
\caption{Contoh format hasil eksplorasi}
\begin{tabularx}{\textwidth}{@{}lXX@{}}
\toprule
Kondisi/Input & Hasil Teoretis & Hasil Pengamatan \\
\midrule
Kondisi 1 & \dots & \dots \\
Kondisi 2 & \dots & \dots \\
Kondisi 3 & \dots & \dots \\
\bottomrule
\end{tabularx}
\end{table}

% ---------------------------------------------------------------------
\subsubsection{Analisis dan Interpretasi}
Hubungkan hasil eksplorasi dengan prinsip kerja dan spesifikasi datasheet.
Jelaskan apakah hasil sesuai teori, penyebab perbedaan, serta implikasinya
terhadap penggunaan komponen pada proyek.

\subsubsection{Keterbatasan dan Sumber Kesalahan}
Bahas noise, toleransi komponen, kesalahan pembacaan, kalibrasi, temperatur,
kondisi lingkungan, keterbatasan library, keterbatasan catu daya, batas
komunikasi, atau faktor lain yang dapat memengaruhi hasil.

\subsubsection{Kesimpulan Pemahaman Komponen}
Tuliskan secara ringkas apa yang dapat disimpulkan mengenai prinsip kerja,
karakteristik sinyal, kebutuhan mikrokontroler, dan kondisi penggunaan
komponen. Tekankan bukti yang menunjukkan bahwa prinsip kerja telah dipahami.


\subsection{Perbandingan Antar-Komponen}
Setelah setiap komponen dibahas secara mandiri, bandingkan karakteristik yang
paling penting untuk menunjukkan alasan pemilihan komponen dan hubungan
masing-masing komponen dengan mikrokontroler.
\begin{table}[H]
\centering
\caption{Perbandingan karakteristik komponen}
\begin{tabularx}{\textwidth}{@{}lXXX@{}}
\toprule
Aspek & Komponen 1 & Komponen 2 & Komponen 3 \\
\midrule
Besaran utama & \dots & \dots & \dots \\
Jenis sinyal & \dots & \dots & \dots \\
Antarmuka MCU & \dots & \dots & \dots \\
Kebutuhan daya & \dots & \dots & \dots \\
Keterbatasan utama & \dots & \dots & \dots \\
\bottomrule
\end{tabularx}
\end{table}

\subsection{Komponen Pendukung}
Jelaskan komponen pendukung yang diperlukan agar komponen utama bekerja
sesuai spesifikasi, misalnya resistor pull-up/pull-down, kapasitor decoupling,
dioda proteksi, transistor driver, regulator, atau level shifter. Hubungkan
setiap komponen pendukung dengan kebutuhan komponen utama yang telah dibahas.
